% P20
\section{Experimental Setup}\label{sec:setup}

\textit{Data.}
We use the SMS Spam Collection~\cite{smsspam}: \num{5574} English SMS messages, \num{747} of them spam.
The collection contains \num{403} exact duplicates of message texts (\num{309} ham, \num{94} spam), typically mass-sent notifications and spam campaigns.
We keep the first occurrence of each text, which leaves \num{5171} messages with \num{653} spam (\SI{12.6}{\percent}).
Without deduplication, \SI{12}{\percent} of the test messages (and \SI{22}{\percent} of the spam among them) would have an identical copy in the training set, which inflates the measured accuracy.

\textit{Split.}
The deduplicated data are split once into a training set of \num{4136} messages (\num{522} spam) and a test set of \num{1035} messages (\num{131} spam), stratified by class with a fixed seed.
All model selection uses the training set only; the test set is used once, for the final results in \cref{sec:results}.

\textit{Model and hyperparameters.}
The plaintext pipeline consists of a TF-IDF vectorizer (lowercased tokens of at least two word characters, L2 normalization), selection of the $k$ features with the highest $\chi^2$ score, and multinomial naive Bayes with additive smoothing~$\alpha$.
TF-IDF and the $\chi^2$ selection are part of the pipeline, so in cross-validation they are refitted on every training fold and do not see the validation fold.
We search a grid of 168 configurations: unigrams or unigrams and bigrams, raw or sublinear term frequency, $k \in \{100, 200, 300, 500, 750, 1000\}$ and $\alpha \in \{0.001, 0.003, 0.01, 0.03, 0.1, 0.3, 1\}$, with stratified 5-fold cross-validation on the training set.
We cap $k$ at \num{1000} because the client's cost grows linearly with $k$.
Since \SI{87}{\percent} of the messages are ham, accuracy separates the models poorly, so we select the configuration with the highest F1 score of the spam class.
The selected model uses unigrams, sublinear term frequency, $k = 1000$ and $\alpha = 0.01$, with a cross-validated accuracy of $0.980 \pm 0.008$ and F1 of $0.915 \pm 0.036$.
Its weights $(b, \mathbf{w})$ are exported as described in \cref{sec:system-design}.

\textit{Baselines.}
We compare the encrypted classifier with three control models evaluated on the same test set:
\begin{itemize}
  \item the \emph{plaintext} naive Bayes model, i.e.\ the same model evaluated in double-precision floating point without encryption or fixed-point encoding; this is the reference for the agreement criteria;
  \item a \emph{majority-class} classifier that always predicts ham, which shows how much accuracy the class imbalance alone gives;
  \item \emph{logistic regression} on the same features, with the regularization strength and class weighting tuned by the same cross-validation ($C = 100$, balanced class weights); it is also a linear model $b + \mathbf{x} \cdot \mathbf{w}$, so it could be evaluated with the same protocol.
\end{itemize}

\textit{Metrics.}
We report accuracy and the precision, recall and F1 score of the spam class.
ROC-AUC is computed from the continuous score: the decrypted log-odds $\tilde{\Delta} / S^2$ for the encrypted classifier and $\Delta$ for the plaintext one.
Agreement is measured as the fraction of test messages with the same label as the plaintext model and the fraction whose decrypted integer equals the plaintext fixed-point computation (\cref{sec:threat-model}).

\textit{Implementation and hardware.}
All experiments run on an Apple M5 Pro (15 cores, \SI{24}{\giga\byte} RAM) with Python~3.14, scikit-learn~1.9 and \texttt{phe}~1.5~\cite{phe} accelerated by \texttt{gmpy2}.
The final results use a \num{2048}-bit modulus $n$ and the scale $S = 10^6$; client-side encryption of the test set is parallelized over 15 processes.
The cost per message is measured separately on a single core, for \num{1024}-, \num{2048}- and \num{3072}-bit keys and $k \in \{500, 750, 1000\}$, averaged over 10 messages.
Random seeds are fixed, so all metrics are reproducible; ciphertexts are not, since \texttt{phe} draws a fresh random $r$ for every encryption.
